\documentclass[11pt,twoside,letterpaper]{article}

% =============================================================================
% 1. CORE SYSTEM PREAMBLE (UNIVERSAL CONFIGURATION)
% =============================================================================
\usepackage{fontspec}           % Modern system font manager
\usepackage{geometry}           % Page margin layout manager
\geometry{letterpaper, margin=1in}
\usepackage{amsmath, amssymb}   % Structural math fonts and symbols
\usepackage{listings}           % Advanced code block container
\usepackage{microtype}          % Fine-tuned typographic spacing adjustments
\usepackage{hyperref}           % Dynamic cross-referencing and hyperlink routing
\usepackage{natbib}             % Advanced academic citation formatting and layout

% --- Safe Font Fallback Configuration ---
\setmainfont{Latin Modern Roman}[
    Extension = .otf,
    UprightFont = lmroman10-regular,
    BoldFont = lmroman10-bold,
    ItalicFont = lmroman10-italic
]

% --- Safe Keywords Command Definition ---
\providecommand{\keywords}[1]{\textbf{\textit{Keywords---}} #1}

% --- Code Listing Aesthetics Setup ---
\lstset{
    basicstyle=\small\ttfamily,
    breaklines=true,
    frame=tb,
    captionpos=b,
    keepspaces=true,
    showstringspaces=false
}

\hypersetup{colorlinks=true, linkcolor=blue, citecolor=blue, urlcolor=teal}

% =============================================================================
% 2. DOCUMENT METADATA INFRASTRUCTURE
% =============================================================================
\title{\textbf{The Contortionist Turn: Diachronic Mapping of the Acro-Yoga Complex from Post-Vedic Statecraft to Late Medieval Somatic Manuals}}

\author{
    \textbf{Patrick S. D. McCartney}\\
    \small Research Affiliate, Anthropological Institute, Nanzan University, Japan \\
    \small School of Culture, History and Language, Australian National University, Australia \\
    \small Centre for Religious Studies, Manipal Academy of Education, India \\
    \small \texttt{u4556787@anu.edu.au}
}
\date{July 29, 2026}

% =============================================================================
% 3. INTRODUCTORY BODY EXECUTION NODE
% =============================================================================
\begin{document}

\maketitle

\begin{abstract}
This paper introduces the \textit{Contortionist Turn} (CT) -- a theoretical and computational digital humanities framework that challenges the dominant, Neo-Romantic historiography of Indian physical culture. Traditional paradigms consistently default to an introspective model, framing postural configurations (\textit{āsana}) and breath-stasis (\textit{stambha}/\textit{prāṇāyāma}) as the organic inventions of isolated forest-dwellers working to transcend material reality. Utilizing an un-truncated, parallel-core semantic search across 20,529 lemmatized text nodes within Hellwig’s \textit{Digital Corpus of Sanskrit} (DCS) and GRETIL repositories, we unearth the volatile, criminalized underbelly of premodern Indian somatic technology. Our sliding-window collocation filters mathematically trace a two-millennium diachronic inversion curve mapping the evolution of the Acro-Yoga Complex across three discrete phases: from post-Vedic statecraft, clandestine espionage (\textit{jambhakavidyā}), and toxicological boundary management (\textit{viṣa-stambhana}) along unstable trade frontiers; through 12th-century courtly and martial capture where dynamic physical shapes were weaponized as close-quarters wrestling submissions (\textit{kūrmāsana}/\textit{niṣpipeṣa}) within state training yards (\textit{akhāḍās}); and culminating in late-medieval monastic sanitization and erasure. We demonstrate that the mechanical components of physical yoga did not grow out of disembodied contemplation; they were the rogue survival tactics of low-status contortionists (\textit{vaṃśa-nartin}, \textit{ḍomba}) and secret field agents systematically captured, regularized, and flattened by elite text-composers into a static, respectably bourgeois display technology (\textit{prekṣaṇīya kām}) engineered to serve contemporary ethno-nationalist soft power and global wellness capitalism.

\vspace{0.3cm}
\keywords{Acro-Yoga Complex, Contortionist Turn, Text-Mining, Statecraft, Subaltern, Espionage, Envenomation}
\end{abstract}

\section{Introduction: De-Spiritualizing the Stasis}
The standard historiography of Indian physical culture consistently defaults to an inward, spiritualized paradigm. Postural configurations (\textit{āsana}) and breathing controls (\textit{prāṇāyāma}) are routinely framed as the introspective inventions of stationary forest ascetics working to transcend the material world. Even recent, critical revisions that unearth the martial origins of pole training (\textit{mallakhāmba}) and its modern adaptation into ``pole yoga'' -- such as McCartney’s foundational work on the \textit{Mānasollāsa} and \textit{Mallapurāṇa} -- remain bounded within elite, courtly, or monastic institutions. This paper introduces the Contortionist Turn (CT): a theoretical and computational framework that strips away this late-medieval theological veneer. By running an un-truncated, parallel-core semantic search across 20,529 digital manuscripts, we expose the material underbelly of Indian somatic technology. Our data reveals that the structural vocabulary of physical yoga operated for a millennium as a volatile network of subaltern survival tactics, clandestine statecraft, and inter-state mercantile transit before it was ever spiritualized by monastic text-composers. 
 
\subsection{The Theoretical Framework: The Materialist Triad and the Guild Network}
To model the historical transmission of these physical techniques without relying on selective qualitative readings, this study conceptualizes the premodern spatial economy through a closed socio-economic infrastructure termed the \textbf{Materialist Triad}. Our digital cross-corpus mapping demonstrates that specialized physical positions (\textit{āsana}) and structural control techniques did not develop out of quiet, un-worldly ascetic isolation. Instead, they emerged within a highly complex, competitive triad consisting of Merchant Guilds (\textit{sārthavāhas}), Acrobat Guilds (\textit{naṭas}/\textit{laṅghakas}), and Yoga Guilds (\textit{akhāḍās}). 
 
These distinct networks interacted continuously along ancient trade routes and frontiers through overlapping matrices of commerce, ritual entertainment, transit security, and state espionage (\textit{gūḍha-caryā}). Within this frontier infrastructure, traveling acrobat-wrestlers operated as liminal, highly adaptable physical specialists. Viewed as otherworldly performers capable of executing acts of ritual awe, their extreme agility, concentration, and fearlessness allowed them to protect commercial transport lines while simultaneously gathering or broadcasting field data. This interaction enabled internal energetic clamps (\textit{bandhas}) to systematically evolve into structural, external postures (\textit{āsanas}). Ultimately, the warrior ascetic yoga guilds institutionalized within regional training yards actively adopted these dynamic, contortionist configurations, weaponizing the physical skills of the subaltern to advance their own socio-political power and secure territorial properties (\textit{yoga-kṣema}).

\subsection{The Vedic Bedrock: Seasonal Martial Mobilization and the Frontier Raid}

The historical baseline of the Contortionist Turn is fundamentally anchored within the geopolitical landscape of the earliest Indo-Aryan frontier. Rather than emerging from passive, introspective forest speculation, the semantic root of \textit{yoga} is structurally bound to seasonal cycles of territorial expansion and violent cross-border conflict. Epigraphical and textual analysis of the Kuru-Pañcāla kings demonstrates an explicit alternation model of statecraft -- forces systematically un-yoked and yoked their resources, setting out on victorious martial raids in the autumn (\textit{yoga}) and returning to fortified settlements for preservation in the summer (\textit{kṣema}). 

In this earliest context, \textit{yoga} defined a literal period of time for martial action, consisting of crossing unstable boundaries, capturing cattle, and dominating regional adversaries. This structural deployment on the frontier required intense, adaptive physical capability to navigate hostile topographies populated by adversarial outgroups. Consequently, \textit{yoga} must be defined historically not as an introverted escape from material reality, but as the systematic application of physical effort and tactical agility engineered specifically for the accumulation of territorial property, symbolic power, and material profit (\textit{yoga-kṣema}). This violent, kinetic substrate provided the structural raw material that elite text-composers and monastic lineages later isolated, sanitized, and turned inward over subsequent millennia.

\subsection{The Subaltern Kinetic Continuum: Pole-Dancing (\textit{Vaṃśa-Nartin}) and Peripatetic Rogues (\textit{Ḍomba})}

The materialist genealogy of the Acro-Yoga Complex achieves definitive historical depth when integrating the ethno-historical tracking of nomadic performer guilds and ritual pole athletics. Deep text layer mapping traces the structural lineage of professional pole play back to the sacrificial matrices of the \textit{Vājasaneīyasaṃhitā} (30.21), where the physical offering of a bamboo pole-dancer (\textit{vaṃśa-nartin}) is explicitly coupled with an outcaste \textit{caṇḍāla} as a sacred dedication to the open air (\textit{antarikṣāya}). Rather than operating as secular side-show amusements, these nomadic specialists deployed extreme bodily contortion, vertical column stasis (\textit{stambha-śrama}), and rope-walking as critical instruments of sympathetic magic engineered to avert collective calamity, ensure agricultural crop regeneration, and re-order cosmic chaos on behalf of the sovereign state.

This subaltern kinetic continuum expands significantly when evaluating the deep text-mining loop for \textit{ḍomba} and \textit{jambhaka} across the narrative layers. In the \textit{Kathāsaritsāgara} (6.2.10420), the \textit{ḍomba} is characteristically represented as an agile, apparatus-bound rogue deploying physical levers, ropes (\textit{pāśa}), and peripatetic performance tools (\textit{mṛdaṅga}) to execute quick physical leaps (\textit{laṅghana}) and escape border checkpoints. Orthoprax legal text-composers responded to this mobile agility with severe institutional anxiety. The structural lineage of physical yoga is thus inexorably bound to this subaltern underbelly -- the \textit{ḍomba-jambhaka} subsystem proves that the mechanical precursors to modern yoga shapes did not emerge from static monastic reflection; they were the rogue survival tactics of low-status contortionists, thieves, and wandering performers navigating an unforgiving geopolitical landscape.

\subsection{The Kautilyan Matrix: Clandestine Extraction and \textit{Jambhakavidyā}}

By the Maurya era, this volatile somatic substrate of the frontier was systematically secularized, institutionalized, and integrated into the centralized espionage operations of the early state. In a critical diagnostic hit from Kautilya’s \textit{Arthaśāstra} (1.12.7), the sovereign’s premier undercover stationary agents (\textit{sattriṇaḥ}) are commanded to master a highly specialized technical curriculum to maintain their performative operational cover along international transit trade lanes:
\begin{quote}
    \textit{ye cāpy asaṃbandhino 'vaśyabhartavyāste lakṣaṇam aṅgavidyām jambhakavidyāṃ māyāgatam āśramadharmaṃ nimittam...}
\end{quote}
Here, the administrative mandate pairs the quantitative study of structural anatomy (\textit{aṅgavidyā}) directly with \textit{jambhakavidyā} -- the subaltern science of metabolic freezing, physical deception, and deceptive camouflage. 

This core interaction reveals that the foundational maneuvers of physical immobilization (\textit{stambhana}) were actively classified as clandestine military and intelligence assets centuries before they were ever spiritualized into monastic configurations. The state did not merely monitor the peripatetic acrobat guilds as peripheral security threats; rather, the administrative elite systematically executed a top-down technological extraction. They vampirized the kinetic and toxicological skills of the subaltern, re-coding the raw mechanics of violent joint-crushing and skeletal compression -- originally derived from unrestricted hand-to-hand combat (\textit{niyuddha}) and predatory clamp-holds (\textit{jambhe dadhmaḥ}) -- into a disciplined, state-sanctioned weapon of counter-surveillance. The \textit{ḍomba-jambhaka} interface is thus exposed as the precise operational valve where rogue survival tactics were formalized by courtly text-composers, setting the template for the long-term scholastic capture and sanitization of the human biological container.

\subsection{The Panoptic Pillar: State Architecture and Confinement}

The earliest high-density intersection isolated on our spatial scatter plot drops squarely into Kautilya’s \textit{Arthaśāstra} (2.5.7). Long before the physical pillar (\textit{stambha}) was celebrated as an asset for martial conditioning or internal yoga alignment, it operated as a literal instrument of state incarceration and panoptic surveillance. The text explicitly mandates the physical parameters required to construct the centralized imperial treasury and state prison complex:
\begin{quote}
    \textit{pakveṣṭakāstambhaṃ catuḥśālam ekadvāram anekasthānatalaṃ vivṛtastambhāpasāram ubhayataḥ paṇyagṛhaṃ koṣṭhāgāraṃ ca dīrghbahuśālaṃ kakṣyāvṛtakuḍyam antaḥ kupyagṛham tad eva bhūmigṛhayuktam āyudhāgāraṃ pṛthagdharmasthīyaṃ mahāmātrīyaṃ vibhaktastrīpuruṣasthānam apasārataḥ suguptakakṣyaṃ bandhanāgāraṃ kārayet} \\
    
    ``He shall cause the state prison (\textit{bandhanāgāram}) to be constructed... featuring columns built of baked bricks (\textit{pakveṣṭakāstambhaṃ}), multiple floor levels and designated stations (\textit{anekasthānatalaṃ}), and open, exposed pillars for surveillance and escape detection (\textit{vivṛtastambhāpasāram})...''
\end{quote}

Kautilya’s administrative blueprint reveals that \textit{stambha} (the pillar) and \textit{sthāna} (the tactical stance/position) were originally engineered as elite architectural tools of spatial and corporeal containment. The specific configuration of the \textit{vivṛta-stambhā-pasāram} -- the open-pillar surveillance grid -- was explicitly designed to ensure that imperial guards maintained total visual mastery over captive bodies, converting the column into a primary apparatus of material restriction. 

The Contortionist Turn argues that the marginalized performer guilds (\textit{naṭas}, \textit{plavakas}) trapped within or monitored by this panoptic architecture systematically weaponized these exact mechanics. They took the physical pillar of state confinement and turned it into a highly skilled, mobile public spectacle of agility and escape -- a subversive performance tradition of structural bypass that was later co-opted and documented in courtly manuals under the sanitized heading of \textit{stambhakrīḍanikā} (pole-play).

\subsection{The Intelligence Nexus: Anatomical Mastery as an Espionage Tool}

Our advanced sliding-window collocation sweep -- forcing the interaction of vital anatomical nodes (\textit{marman}) and breath pathways (\textit{prāṇa}) alongside clandestinity (\textit{gūḍha}) and performance (\textit{naṭa}) -- pierces directly into the tactical reality behind these somatic skills. Kullūkabhaṭṭa’s classic medieval commentary on the \textit{Manusmṛti} (Verse 7.154) explicitly links high-level anatomical knowledge directly to state espionage:
\begin{quote}
    \textit{kāpaṭikodāsthitagṛhapatika-vaidehikavāpasaṃvyañjanātmakaṃ pañcavidhaṃ caravargaṃ tattvatacintayet | tatra kāpaṭikaḥ paramarmajñaḥ pragalbhaś chātraḥ...}
\end{quote}

The state’s premier undercover field agent (\textit{kāpaṭika} / \textit{cara}) is thus lexicographically defined by two explicit attributes: absolute mastery over fraud, disguise, and occupational infiltration (\textit{kapaṭa-vyavahāra}), and a precise, lethal expertise in diagnosing and exploiting the vulnerable, vital nodes of an opponent's body (\textit{paramarmajña}). 

To ground this operational intersection, the Contortionist Turn rejects the assumption that the state handed corporate espionage badges directly to outcaste acrobats; rather, the centralizing administrative elite executed a top-down technological extraction. The state vampirized the kinetic agility and anatomical precision developed within peripheral performer networks, re-coding the material realities of aggressive, hand-to-hand combat (\textit{niyuddha}) into a weaponized intelligence asset. In the epic accounts of arena duels preserved in the \textit{Viṣṇu Purāṇa} (5.20.66), physical configurations (\textit{āsana}) are contextualized not as instruments of internal isolation, but as violent mechanisms of skeletal destruction:
\begin{quote}
    \textit{prāharan muṣṭinā mūrdhni vakṣasyāhatya jānunā} \\
    \textit{pātayitvā dharāpṛṣṭhe niṣpipeṣa gatāyuṣam}
\end{quote}
Here, the physical manipulation of the joints -- executing an arm-clapping challenge, kneeing the chest, and grinding an opponent into the earth (\textit{niṣpipeṣa}) -- demonstrates that specialized bodily training (\textit{mallapariśrama}) belonged to a violent, combat-ready continuum long before it was theoretically cleaned up by monastic text-composers. Secret agents disguised themselves as wandering actors or ascetics (\textit{chadmanā}) to slip past border checkpoints precisely because the state recognized this subaltern physical mastery as the supreme asset for infiltration and covert execution.

\subsection{The Geopolitical Formula: From Vedic Trampling to Arthaśāstra Utthāna}

The operational transition of \textit{yoga-kṣema} from an epic, pastoral binary into a strategic technology of statecraft is explicitly codified by pairing late-Rigvedic conquest metrics with Kautilyan political science. In the ritual architecture of \textit{Ṛgveda} (7.54.3), the binary division of communal activity is strictly partitioned between peace and active warfare (\textit{pāhi kṣema uta yoge}). This framework is inherently aggressive; as detailed in \textit{Ṛgveda} (10.166.5), the successful seizure of an adversary’s mobilization capacity and domestic security (\textit{yoga-kṣemaṃ va ādāyāham bhūyāsam}) is structurally validated through visceral, bodily domination and trampling (\textit{uttama ā vo mūrdhānaṃ akramim}). 

By the period of the classical Maurya state, Kautilya’s \textit{Arthaśāstra} (1.7.1) completely secularizes this alternation model into a precise programmatic mandate for the sovereign, commanding the execution of political goals through relentless, active physical effort (\textit{utthānena yoga-kṣema sādhanam}). Within this Kautilyan framework, \textit{yoga} defines the successful, pragmatic accomplishment of state objectives -- specifically the acquisition of unpossessed resources -- while \textit{kṣema} defines the preservation and peaceful enjoyment of that acquired welfare. Specialized physical disciplines (\textit{utthāna}) were thus engineered to service the expansionist requirements of the centralized state long before they were adapted into introverted monastic physical regimes. The dynamic-static continuum of \textit{yoga-kṣema} underscores that state preservation is routinely secured through forceful, non-peaceful boundary management.

\section{Methodology: The Multi-Source Grid and Sliding-Window Collocation}

To track these structural and semantic shifts across two millennia without falling into phonetic traps, this study implements a self-contained, parallel-core Python processing pipeline. Traditional raw string scanning fails when Sanskrit words encounter complex phonetic mutations (\textit{Sandhi}), rendering standard keyword lookups scientifically inaccurate. To bypass this lexical instability, our engine integrates Oliver Hellwig’s \textit{Digital Corpus of Sanskrit} (DCS), which systematically breaks text nodes down into individual, lemmatized data columns, mapping every word back to its dictionary base form (\textit{lemma}) regardless of contextual spelling changes.

Crucially, our macro-pipeline does not merely calculate abstract keyword densities -- an approach that cannot automatically deduce localized operational nuance. Instead, the computational engine operates as a highly calibrated \textbf{sliding-window collocation filter}. By configuring a fixed token window ($\pm 5$ words), the script suppresses ambient semantic noise and programmatically isolates high-density syntactic nodes where specialized somatic terms (\textit{āsana}, \textit{stambha}) co-occur alongside explicit variables of statecraft, toxicology, or subaltern performance. This methodology mathematically validates our subsequent close-readings, isolating genuine text-historical smoking guns from arbitrary textual baseline noise.

Our script evaluates five primary thematic categories across a merged corpus of 20,529 text nodes:
\begin{lstlisting}[language=Python, caption={Computational Lexicon Mapping Configurations}, label={list:lexicon}]
SOMATIC_LEXICON = {
    "postural_contortion": ["āsana", "pīṭha", "vṛścika", "padma", "śalabha"],
    "apparatus_pole": ["vaṃśa", "stambha", "stamba", "veṇu"],
    "subaltern_tribal": ["caṇḍāla", "ḍomba", "plavaka", "jambhaka", "ābhīra"],
    "poison_necromancy": ["viṣa", "gara", "śmaśāna", "bhūta", "śava"],
    "espionage_subversion": ["gūḍhapuruṣa", "cara", "sattrin", "kāpaṭika", "chadman"]
}
\end{lstlisting}
By aggregating these outputs into two composite scores -- the Subversive Mobility Index (Transit + Espionage + Leaping) and the Somatic Contortion Index (Postures + Poles) -- our engine isolates the definitive non-linear mathematical outliers in the dataset. This quantitative framework maps a shared stealth vocabulary running between military scouts, clandestine border agents (\textit{cara}), and traveling acrobat guilds (\textit{naṭas}), where the physiological capacity to compress, twist, and glide the body (\textit{saṃsarpaṇa} / \textit{upasarpaṇa}) operated as a raw material mechanic of frontier survival.

\subsection{Diachronic Spatial Analogues: The Ethnographic Mapping of the Naṭ and Kabūtari}

The historical trajectory of the subaltern performer guild (\textit{naṭa}) is further contextualized by evaluating modern ethnographic records from the colonial frontier. Twentieth-century registers, such as Russell’s (1916) survey of the Central Provinces, track a highly mobile complex of marginalized outgroups, including the Bādi, Dang-Charha (rope/pole climbers), Bāzigar, and Sapera (snake-charmers). To maintain methodic unity and avoid the anachronistic trap of assuming an organic, primordial continuum over two millennia, this study frames these modern ethnographic records strictly as an \textbf{analogue of structural positioning} rather than a linear, biological lineage. 

The material conditions of peripheral peripatetic groups -- where advanced bodily contortion, border-evasion, and physical plasticity serve as core economic currency -- remain structurally constant across historical epochs due to the stable exclusionary mechanics of the orthoprax state. For example, female acrobats within these guilds are labeled as \textit{kabūtari} (pigeons) due to their execution of extreme aerial backbends and tumbling drops that mimic the specialized flight loops of the tumbler pigeon (\textit{Columba gyratrix L.}). Rather than an isolated survival, this performance matrix represents a structural analogue to post-Vedic definitions of shamanic flight (\textit{ut-mad}) and monastic classifications of pole-vaulting (\textit{laṅghanakammaṃ}). The tight social overlap between the pole-climbing Dang-Charha and the toxicological Sapera demonstrates that vertical column training (\textit{stambha-śrama}) and venom-freezing (\textit{viṣa-stambhana}) operated as part of a single, shared pool of frontier physical technologies maintained by itinerant communities centuries before they were sanitized into monastic hatha configurations.

\subsection{The Enclosed Column: Pylwan Gopauls, Kolhāṭi Poles, and the State Surveillance Grid}

The structural link between mobile combative agility and vertical column conditioning is definitively verified by historical surveillance records tracking the outgroups of the Deccan interior. Nineteenth-century administrative ledgers, such as Gunthorpe’s (1882) \textit{Notes on Criminal Tribes}, document that itinerant communities designated as the Pylwan Gopauls (Wrestler Gopals) operated as a nomadic, marginalized population whose primary livelihood depended on a combination of pugilistic wrestling (\textit{pehlavān}) and the execution of high-level gymnastics performed on a long mobile pole (\textit{khām}). This apparatus-bound mobility is reinforced by the \textit{Census of India 1901} (Volume VIII: Berar), which traces the etymology of the wandering Kolhāṭi acrobat guild directly to the word \textit{kolah} -- the long bamboo pole utilized explicitly for jumping, vaulting, and physical leaping (\textit{laṅghana}). 

The census registers confirm that these pole-play lineages, alongside parallel sub-castes like the Bānsberias (bamboo-performers) and Jaithis (\textit{Jyeṣṭhi-mallas}), shared an identical subaltern social stratum characterized by rope-dancing, public entertainment, and systemic social exclusion. By utilizing this structural-analogue framework, the Contortionist Turn maps a robust somatic genealogy running from ancient tribal outgroups (\textit{caṇḍālas}) to medieval arena specialists, and finally to modern criminalized clusters. The physical technology of the pole did not emerge from stationary monastic reflection; it was a mobile, functional asset of traveling guilds packed onto draft animals across the Deccan, which elite text-composers systematically captured, institutionalized, and stripped of its subaltern identity to serve an orthodox theological narrative.

\subsection{The Indigenous Substrate: Kol Cosmologies and Shamanic Pole-Ascent}

The diachronic genealogy of the Acro-Yoga Complex is deeply rooted within an indigenous, Austroasiatic and Dravidian tribal substrate that predates courtly and monastic institutionalization. Ethnographic and spatial mapping of the Munda-speaking Kols (including the Santāls) exposes a foundational somatic environment where vertical column stasis (\textit{stambha}) operated as a shamanic ritual technology rather than a static physical exercise. Within Kol cosmological frameworks, the physical act of scaling high vertical bamboo poles functions as a literal mechanism of ritual ascent to commune with cosmic polarities, specifically the Sun god (\textit{Siṅ Boṅga}) and the Moon goddess (\textit{Ñinda Chando}), alongside more complex Dravidian configurations of the Great Mother Goddess and Sky Father. 

This indigenous practice establishes the material raw material of pole conditioning within a non-Brahminized, forest-dwelling topography. Geographically concentrated across the Kolha macro-corridors of the subcontinent, these tribal populations represent the direct historical ancestors of the peripatetic acrobat guilds documented in later colonial tracking registers as the pole-leaping Kolhāṭis and classical Kollāṭas. The trajectory remains clear: the somatic mechanics of vertical apparatus athletics were systematically extracted from the ritualized, ecstatic climbing traditions of indigenous Adivasi communities before being captured by regional states for fortification science, and ultimately sanitized into the introspective postural technologies of late-medieval hatha yoga.

\section{Applied Yoga and the Geopolitics of Erasure: From Kautilyan Matrices to Modern Yogaland}

The historical trajectory of the Acro-Yoga Complex achieves its ultimate ideological encapsulation within the transnational consumption-scape of modern ``Yogaland'' \cite{McCartney2021}. In contemporary state-led branding initiatives and global wellness forums, the concept of ``Applied Yoga'' is structurally deployed as a Neo-Romantic panacea for environmental degradation, climate crises, and social instability \cite{McCartney2021}. This policy proposal rests upon a pervasive article of faith: that ancient forest-dwelling ascetics operated as nature-loving, proto-environmentalist figures who sourced their soteriological insights from peaceful biocentric communion \cite{McCartney2021}. A rigorous, diachronic interrogation of primary Sanskrit sources completely disconfirms this disembodied myth \cite{Birch2024, McCartney2021}. Close reading reveals that long before physical techniques were turned inward as static, spiritualized technologies of self-cultivation, they functioned as strategic instruments of settler-colonial expansion, border fortification, and territorial acquisition (\textit{yoga-kṣema}) at the volatile frontier of the early state \cite{Birch2024, McCartney2021}.

\subsection{The Bio-Moral Monopoly: Greenwashing the Frontier}

This structural amnesia underpins the modern alignment of yoga-inflected embodiment with bio-ecological popularism and majoritarian nationalism \cite{McCartney2021}. As the contemporary state leverages the concept of \textit{sanātana dharma} to assert an exceptionalist model of preventative healthcare and soft power, it executes a calculated horizontal expansion across the global medical tourism marketplace \cite{McCartney2021}. This greenwashed scaffolding deliberately backgrounds the volatile, material history of its constituent parts \cite{McCartney2021}. 

Modern yoga nationalism strips away the nomadic, criminalized contexts of the subaltern performer guilds (\textit{naṭas}, \textit{laṅghakas}), outcaste root-cutters (\textit{śvapāka}), and clandestine field intelligence purveyors (\textit{gūḍhapuruṣas}) \cite{McCartney2021}. By re-coding dynamic, dangerous assets of frontier survival into a highly managed, clock-bound daily practice schedule, the rogue body is successfully domesticated \cite{McCartney2021}. The transformation is finalized when the grueling mechanics of manual combat holds are reduced to a spectacular display technology (\textit{prekṣaṇīya kām}) designed for public consumption -- proving that whether turned outward as an instrument of statecraft or inward as passive devotion, the mechanical vocabulary of the body is systematically co-opted by elite text-composers to enforce institutional power, regulate human biological assets, and protect territorial prosperity \cite{McCartney2021}.

\subsection{The Spatial Construct: Deconstructing the Myth of the 'Vedic Village'}

The architectural and ideological scaffolding of this civilizational erasure achieves its structural localization within the systematic deconstruction of the mythic 'Vedic Village' \cite{McCartney2026b}. Rather than validating the historical village as an introspective, nature-centered ecological haven, digital cross-corpus mapping exposes the rural layout as a deeply contested symbolic and political artifact engineered for top-down settler-colonial expansion \cite{McCartney2026b}. The primitive village functioned as a strategic frontier anchor under the alternation mechanics of \textit{yoga-kṣema} -- enforcing strict spatial, ritual, and caste boundaries to isolate corporate wealth and advance majoritarian state consolidation \cite{McCartney2026b}.

This exclusionary spatial enclosure was maintained through explicit somatic violence designed to patrol access to the elite linguistic and ritual core \cite{McCartney2026b}. As documented in the punitive constraints of the \textit{Gautamadharmasūtra} (12.4--6), the preservation of liturgical purity required visceral corporeal mutilation, explicitly commanding the pouring of molten tin (\textit{trapu}) into the ears of the unauthorized low-caste listener and the dismemberment (\textit{bheda}) of the body for structural retention (\textit{dhāraṇa}) of the text \cite{McCartney2026b}. The 'Vedic village' is thus exposed by the Contortionist Turn not as a historical demographic reality, but as an aspirational construct fueled by modern ethno-nationalist nostalgia to greenwash a highly stratified and militarized past \cite{McCartney2026b}. By tracing the overlapping matrices of merchant guilds (\textit{sārthavāhas}) and mobile acrobat spies (\textit{naṭas}), our pipeline confirms that the physical and linguistic canons of modern yoga were systematically extracted from this volatile, subaltern frontier underbelly before their late-medieval monastic sanitization \cite{McCartney2026b}.

\section{The Weaponized Breath: Paralyzing the Frontier}

This toxicological and martial intersection is reinforced by our text-mining hits within the \textit{Kauśika Paddhati}, the premier practical and ritual handbook of the Atharvaveda \cite{KauśP}. Long before medieval hatha texts re-coded bodily stasis (\textit{stambha}) as an internal, contemplative breath-retention hold (\textit{kumbhaka}), early semantic layers openly categorized it as a weaponized technique of physical immobilization, combat subversion, and chemical/alchemical extraction \cite{Birch2024, KauśP}:
\begin{lstlisting}[language=Sanskrit, caption={Kauśika Paddhati Overlap Configuration}, label={list:kaushika}]
mohanaṃ stambhanamityarthaḥ | yairamurchāacetanāṃ sukhaṃ hanyante | 
śāntyāṃ hastyavapadātīnāṃ sarveṣāṃ mohanamacetanatvaṃ bhavattyarthaśāstrakāraiḥ...
...bhaiṣajyāni āstambhnadgrathnanti ... viṣastambhanabhaiṣajyam ... 
yena viṣaṃ na visarpati | śarīreasthitaṃ bhavati
\end{lstlisting}

The text explicitly equates military sorcery and tactical subversion (\textit{stambhanam}) with inducing catatonia and absolute unconsciousness (\textit{acetanatvam}) in opponents so they can be easily neutralized on the battlefield \cite{KauśP}. Crucially, this overlaps with \textit{viṣa-stambhana} -- the literal binding, freezing, or arresting of poison inside the body using localized somatic locking techniques (\textit{granthi}) to prevent systemic venom circulation (\textit{na visarpati}) \cite{KauśP}. Controlling the breath and locking the limbs operated first as raw physiological survival mechanisms for traveling performers, snake-handlers (\textit{gāruḍas}), and shadow agents navigating toxic or hostile frontier trade routes \cite{Birch2024}. To build the essential intermediate step in our transmission model, we argue that medieval ascetic orders straddled the boundary between these wandering forest root-cutters and institutionalized monastic houses, inheriting the physical need to survive toxic topographies before turning that vocabulary of physical containment inward (\textit{stambhakarī mudrā}) to serve a new metaphysics of internal fluid retention (\textit{bindu-dhāraṇa}) \cite{Birch2024}.

\subsection{The Epistemology of Verticality: State Surveillance and the Normalization of the Order}

The institutional mechanism of the Contortionist Turn receives absolute structural validation when read against classical and contemporary histories of ascetic containment \cite{Clark2006, Olivelle1987}. Rather than developing within an unchanging spiritual vacuum, our pipeline confirms that early physical behaviors were aggressively monitored and controlled by the centralizing state \cite{Olivelle1987, Olivelle2011}. As documented by Patrick Olivelle, classical administrative policy under the Kautilyan matrix did not treat the ascetic self as a transcendent entity, but as a dangerous vector of fluid border mobility requiring strict political regulation and state-sanctioned enclosure \cite{Olivelle1987, Olivelle2011}. This structural containment reflects the exact historical moment where the volatile, independent skills of outcaste contortionists and field spies were captured, alphabetized, and compressed under the corporate order of the \textit{Daśanāmī Saṃnyāsīs} to serve as disciplined military labor for regional courts \cite{Clark2006, Clark2018}.

This historical process of normalization directly explains the dramatic vertical spikes tracked along our diachronic models \cite{McCartney2026b}. As analyzed by Chris Danta through the lens of vertical aesthetics, the evolution of physical yoga marks a calculated transition from the rogue agility of the peripatetic acrobat to the rigid, static discipline of the institutionalized ascetic \cite{Danta2015}. The grueling mechanics of physical locking (\textit{stambhana}) and manual close-quarters combat submission holds were stripped of their nomadic mobility and repurposed as spectacular display technologies (\textit{prekṣaṇīya kām}) designed to project state civilizational authority to civic audiences \cite{Danta2015, McCartney2026b}. The multi-millennial transformation is finalized when the dangerous, rogue skills of the frontier underbelly are successfully turned inward and moralized, transforming the body into a highly aestheticized, bourgeois instrument of passive devotion and transnational capitalist wellness tourism \cite{McCartney2021, Srinivas2012}.

\subsection{The Topology of Somatic Infiltration: Quantitative Extraction and Scholastic Capture}

The final comparative processing loop for the subaltern lemmas \textit{ḍomba} and \textit{jambhaka} isolates exactly 248 high-density syntactic nodes, providing the empirical foundation required to ground our diachronic models within concrete histories of military labor, institutional containment, and subaltern evasion \cite{Lorenzen1978, McCartney2026c}. In the earliest chronological layers, the root \textit{jambha} functions as a literal and ritual engine of physical crushing, preserved across the parallel core nodes of the \textit{Maitrāyaṇīsaṃhitā} and the \textit{Atharvaveda} which command the forceful binding of an adversary within a predatory skeletal hold: \textit{yo 'smān dveṣṭi yaṃ vayaṃ dviṣmas taṃ vo jambhe dadhmaḥ} \cite{McCartney2026c}. This somatic archive, read against Rosalind O'Hanlon's mapping of martial culture, demonstrates that long before late-medieval text-composers re-coded locking mechanisms (\textit{bandha}) as internal contemplative seals, the technical vocabulary of physical constraint relied entirely on manual joint-crushing and skeletal clamping derived from wrestling pit gymnasia (\textit{tālīm khānā}) and close-quarters grappling (\textit{malla-yuddha}) \cite{Birch2024, OHanlon2007}.

By the Mauryan era, this volatile somatic substrate was systematically secularized and integrated into centralized espionage operations \cite{McCartney2026b}. As codified in the critical diagnostic hit of \textit{Arthaśāstra} (1.12), the state's premier undercover stationary agents (\textit{sattriṇaḥ}) were commanded to master a technical curriculum that explicitly coupled structural anatomy (\textit{aṅgavidyā}) directly with \textit{jambhakavidyā} -- the subaltern science of metabolic freezing, physical deception, and deceptive bodily camouflage \cite{McCartney2026b}. This tactical capability mirrors the material survival mechanics of the criminalized \textit{ḍomba} guilds tracked across the narrative layers of the \textit{Kathāsaritsāgara} (\textit{KSS 6.2}: \textit{mṛdaṅgahasto moṣāya ḍombaḥ}) \cite{McCartney2026c}. Characterized by their reliance on apparatus-bound performance levers, ropes (\textit{pāśa}), and peripatetic mobility, these traveling contortionists deployed extreme physical agility as a dual asset for public spectacle and clandestine border-evasion \cite{McCartney2026c, McCartney2026b}. 

As Patrick Olivelle demonstrates, the premodern state aggressively monitored this fluid mobility, leading to a calculated historical capture where regional courts systematically co-opted these rogue frontier assets and compressed them into the regimented corporate orders of the militant \textit{akhāḍā} \cite{Clark2006, Olivelle1987}. As established in the foundational historiography of William Pinch, these peripatetic networks operated as premier armed mercenary cartels and trans-regional trading syndicates whose ascetic identity functioned as an absolute corporate camouflage for long-distance intelligence gathering and illicit capital transit across highly contested regional thresholds \cite{Pinch1996, Pinch2012}. The evolution of physical yoga is thus shown to be structurally parasitic upon this subaltern underbelly; elite text-composers captured these dynamic combat and defense skills and flattened them into the sanitized, respectably bourgeois spectacular display technologies (\textit{prekṣaṇīya kām}) of modern postural catalogs \cite{Birch2024, Danta2015}.

\subsection{Tribal Topographies and Necromantic Camouflage on the Frontier}

The manual collection of data from the \textit{Harivaṃśa} appendices explicitly anchors early combative and physical agility within an outcaste, frontier landscape rather than a static monastic environment \cite{McCartney2026b}. Appendix 1 maps the cult of the Goddess (\textit{Mahādevī}) directly onto un-Brahminized forest regions, noting her structural worship by marginalized tribal networks:
\begin{quote}
    \textit{parvatāgreṣu ghoreṣu nadīṣu ca guhāsu ca /} \\
    \textit{śabarair barbaraiś caiva pulindaiś ca supūjitā //}
\end{quote}
This hazardous environment of extreme terrain and continuous, transgressive movement indicates that specialized physical agility was initially developed as an adaptive survival trait by nomadic forest-dwellers like the Śabaras and Pulindas.

This frontier topography intersects directly with toxicological and mortuary environments, providing the concrete setting for what our framework defines as necromantic camouflage. The epic text documents a perilous space populated by raw-flesh eating entities (\textit{kravyāda-gaṇa}) and flesh-eating ghouls (\textit{piśāca}) moving toward urban administrative centers like Karavīrapura -- a site explicitly bound to arsenic production, cemetery fields, and sorcery spells. To survive and operate within this terrain, clandestine state agents (\textit{cara}) and traveling performer guilds (\textit{naṭas}) deployed specialized bodily contortions. They used the frantic, rhythmic dances of flesh-eating attendants (\textit{ḍākinyo ’tra pranṛtyanti}) or the rigid, corpse-like immobility of the charnel ground (\textit{stambha}) as a deceptive mechanism to slip past highway checkpoints undetected. When these performer networks were mobilized during mass state festivities (\textit{mahāghoṣa}), their public displays of tumbling and physical agility provided the ultimate espionage cover for secret agents (\textit{gūḍhapuruṣa}) trained in targeting the body’s vital channels (\textit{marman}).

\subsection{The Frontier Roots: Śvapāka Root-Cutters and the Kalpasthāna Grid}

The structural link between subaltern mobility and the toxicological frontier is further verified by tracking the occupational alternatives of outcaste communities across the classical technical lexicons. While orthoprax legal digests like the \textit{Manusmṛti} systematically restrict the degraded \textit{śvapāka} or \textit{sopāka} lineages to the defiled office of public executioner, alternative lexical extractions explicitly document their practical role as physical specialists who collect and sell medicinal roots. This frontier botany is structurally bound to the manipulation of lethal chemistry and sorcery -- the lemma \textit{asita} maps interchangeably onto the night, dark black snakes, the planet Saturn, and the explicit lord of darkness and magic in the \textit{Atharvaveda}. 

This marginalized monopoly over potent mineral vitriol (\textit{tuttha}) and toxic plant formulations (\textit{mahāpañcaviṣa}) was institutionalized within traditional medical frameworks under the specific heading of the \textit{kalpasthāna} -- the specialized science of poisons and antidotes. Wandering acrobat-grapplers and snake-handlers (\textit{saperas}) utilized deep breath retention and structural physical locks (\textit{viṣa-stambhana}) as critical physiological levers to control internal vital energy (\textit{prāṇa}) and arrest the spread of systemic toxins while crossing unstable regional boundaries. The Contortionist Turn thus demonstrates that the physical and linguistic components of modern yoga were systematically extracted from this volatile, weaponized arena of veterinary medicine, frontier defense, and state intelligence gathering before their late-medieval theological sanitization.

\section{The Courtly Capture: Postural Weaponization and Caste Shifting}

By the 12th century, our diachronic timeline maps a major structural shift. The raw, volatile, border-crossing body skills of the subaltern were systematically captured and institutionalized by elite regional courts. King Someśvara III’s \textit{Mānasollāsa} (\textit{Mallavinoda} section) captures this transition perfectly, documenting the moment elite martial arts absorbed these physical configurations:
\begin{quote}
    \textit{mūrdhni kūrmāsanaṃ badhdvā bhajedvāpyekapādake /} \\
    \textit{uttānapratimallasya nivārya caraṇadvayam //}
\end{quote}
The text commands a courtly wrestler (\textit{malla}) to actively bind and lock the tortoise posture (\textit{kūrmāsanam badhdvā}) directly on the head or neck of their opponent (\textit{mūrdhni}) to pin them down. 

This exposure completely subverts the modern ascetic myth. Advanced non-seated positions like \textit{kūrmāsana} were not invented by peaceful forest meditators; they were aggressive, weaponized grappling submissions used in regional wrestling pits (\textit{akkhāḍaka} / modern Akhāḍā). High-caste patrons and text-composers took the physical skills of traveling acrobat guilds, stripped away their nomadic context, and re-coded them into an elite, courtly martial system. This process allowed specialized performer lineages (like the Jyeṣṭhimallas or Modha Brahmins) to successfully climb the social ladder and claim high-caste status while retaining their monopoly over physical culture.

\subsection{The Biomechanical Mechanics of Niyuddha and Pachāṛ}

The manual collection of technical martial terms demonstrates that physical postures (\textit{āsana}) and structural locks (\textit{stambha}) were originally rooted in the violent realities of bone-crushing combat (\textit{niyuddha}). Vernacular Indo-Persian wrestling registers document precise vectors of skeletal destruction. The \textit{gā’o pachāṛ} (or \textit{qasāī’ā dāw}) handles structural leverage by using a high neck lever to force an opponent into systemic biomechanical collapse, mimicking the slaughterhouse mechanics of forcing down large cattle. This low-line tactical collapse is reinforced by targeting the \textit{agrajaṅghā} (the shin-bone) or \textit{pūrvasaktha} (upper thigh) during a \textit{ghoṛā pachāṛ} strike to break the opponent’s foundation.

This micro-anatomical targeting aligns perfectly with Epic accounts of unrestricted arena combat. In the \textit{Harivaṃśa} (75.42), Kṛṣṇa’s final strike against Caṇūra is defined as an un-weaponized execution (\textit{aśastram atighoraṃ}) consisting of raw physical power (\textit{balaprāṇaviniṣpādyaṃ}):
\begin{quote}
    \textit{prāharan muṣṭinā mūrdhni vakṣasyāhatya jānunā} \\
    \textit{niḥsṛte sāśrurudhire tasya netre sabandhane}
\end{quote}
The subsequent execution of a crushing ground-hold (\textit{niṣpipeṣa}) in the \textit{Viṣṇu Purāṇa} (5.20.66) demonstrates that deep floor holds were originally weaponized grappling submissions designed to stop an opponent’s vital breath (\textit{prāṇa}). Furthermore, the directional stealth vocabulary of \textit{saṃsarpaṇa} (sneaking, gliding, or launching surprise attacks) and \textit{upasarpaṇa} (deceptive creeping) bridges the tactical space between martial combat and the border-crossing agility of low-caste acrobat guilds (\textit{naṭas}).

\subsection{The Acoustic Arena: Yogakṣema and the Alternation Model of Combat}

The linguistic foundation of the Contortionist Turn is definitively anchored by tracking the terms \textit{yoga} and \textit{kṣema} back to their earliest attested layers in the Indo-Aryan corpus. As codified in \textit{Mallapurāṇa} (7.1), the acoustic environments and kinetic shouts performed within the training yard are explicitly systematized for material acquisition and political dominance:
\begin{quote}
    \textit{ākhāḍhasya samīpasthā ye kāryāḥ svarāḥ kila |} \\
    \textit{tatsarvaṃ kathayiṣyāmi yogakṣemārthasiddhaye ||}
\end{quote}
This verse demonstrates that the sounds generated around the wrestling pit (\textit{ākhāḍha}) were engineered strictly for the realization of material security (\textit{yogakṣemārthasiddhaye}). 

This martial register directly preserves the earliest semantic structures of the \textit{Ṛgveda} (10.166.5), where \textit{yoga} defines alternating periods of mobilization or martial action, and \textit{kṣema} defines periods of non-martial rest or structural preservation. Within the technical lexicon of the \textit{Mallapurāṇa} (13.34), \textit{yoga} is explicitly equated with the mechanics of wrestling combat (\textit{malla-yuddha}), and physical postures (\textit{āsanas}) are defined as active ``fighting postures'' or ``modes of fighting.'' The acquisition of physical mastery through rigorous exertion (\textit{jita-śrama}) was deployed with the ultimate territorial intention of total conquest (\textit{trailokya-vijayī}). This structural lineage tracks a clear historical path: long before late-medieval hatha yoga manuals re-coded \textit{āsana} as a static, internal tool for spiritual liberation, the configurations of the body operated as kinetic instruments of physical dominance designed to secure material prosperity and defend regional state boundaries.

\subsection{The Dynamic-Static Binary: Debunking the Upanisadic Union Myth}

The conceptual mechanics of the Contortionist Turn require a rigorous deconstruction of popular, ahistorical tracking models that default to an immediate, mystical definition of yoga as spiritual ``union.'' Textual mapping of the earliest Vedic strata demonstrates that the social, ritual, and political world of Indo-Aryan society pivoted fundamentally on the binary matrix of \textit{yoga} and \textit{kṣema}. Operating as explicit indicators of time, \textit{yoga} denoted the period for martial action -- specifically the material act of yoking the body to weapons, armor, and chariots for cross-border raids -- while \textit{kṣema} denoted the corresponding period for rest, security, and the static preservation of wealth. 

The totality of \textit{yoga-kṣema} must therefore be understood as a native dynamic-static continuum. While modern popular paradigms default to later Upanisadic re-codings (such as the \textit{Kaṭha Upaniṣad}) to frame breath restraint and upright stasis as detached tools to ``overcome death,'' a sociological audit reveals that even these internalizations operated as tactical actions engineered for the acquisition of power, otherworldly prosperity, and protective symbolic armor. Long before late-medieval hatha text-composers isolated individual bodily contortions as spiritual gymnastics, the physical configurations of the organism were functionally deployed within this alternation framework to master material transitions between combative mobilization and structural grounding.

\subsection{The Kinetic Continuum: Lāga, Karaṇa, and the Performance Spectrum}

The technical taxonomy preserved in Chapter 8 of the \textit{Mallapurāṇa} completely de-spiritualizes the early mechanics of posture, mapping \textit{āsana} and \textit{sthāna} as components of an active, kinetic continuum rather than static monastic configurations. The text outlines a highly structured catalog of fourteen mechanical grips known as \textit{Lāga} (Holds) spanning verses 5a--9b, dictating explicit vectors of close-quarters anatomical binding, including hair-seizing, shoulder-levers, backbone locks, and supportless abdominal compressions. 

Crucially, this combative movement array is structurally linked to classical performance and dance tracking sciences. As codified across the \textit{Nāṭyaśāstra} (4.56, 4.61) and the \textit{Saṅgītaratnākara} (4.7), complex bodily actions (\textit{karaṇas}) are constructed by stacking movements from preceding structural layers, requiring a practitioner to master the split-second transition between being stationary (\textit{sthāna}) and moving (\textit{cārī}) to achieve coordinated kinetic fluidness (\textit{aṅgahāra}). This shared somatic substrate operating interchangeably between wrestlers, dancers, and acrobats is executed via a strict physical hierarchy: prioritizing lower-frame footwork (\textit{calapāda}), arm clearing, and micro-adjustments of the palms (\textit{karatalayoścalana}). This combative attitude of readiness to secure an absolute footing survives directly within modern regional physical registers as the Marathi wrestling term \textit{pavitrā} \cite{Patwardhan1927}, verifying that medieval martial conditioning mechanisms provided the technical raw material that late-medieval text-composers later isolated and turned inward.

\subsection{The Spatial Corridor: The Jyeṣṭhi-malla Migration and the Kalyani Fort Nexus}

The institutionalization of these physical techniques within the 12th-century Chalukyan court is fundamentally bound to a massive, state-sponsored spatial migration corridor known as ``The Wrestlers’ Route.'' Historical and epigraphical records confirm that c. 1125--1129 CE, approximately 350 Jyeṣṭhi-malla families migrated south from Modhera (Mehsana District, Gujarat), traversing the hazardous, un-Brahminized frontiers of the Vindhya and Satpura mountain ranges to settle directly within the Western Chalukyan capital of Kalyani Fort (Basavakalyan, Karnataka). This migration trajectory grounds our digital text-mining metrics in a concrete, militarized landscape. 

Architectural and defensive surveys of the Kalyani Fort complex isolate a highly fortified, subterranean chamber designated as the \textit{Tālīm Khānā} -- a dedicated military garrison gymnasium built directly into the fortress walls and surrounded by defensive ditches. This spatial and architectural configuration demonstrates that specialized postural development was structurally integrated into royal fortification science and state-controlled training yards (\textit{akkhāḍaka}) long before it was adapted into contemplative monastic environments. The Jyeṣṭhi-mallas did not transport a disembodied spiritual philosophy; they transported a highly functional, weaponized system of close-quarters wrestling drills (\textit{mallapariśrama}) and vertical column conditioning (\textit{stambha-śrama}) designed to sustain the defensive infrastructure of the regional state.

\section{The Diachronic Climax: Theological Cleaning and Static Isolation}

The final chronological phase of our diachronic pipeline captures the ultimate structural inversion of the Acro-Yoga Complex. As the independent, fluid mobility of the frontier underbelly crashed downward between the 15th and 18th centuries, late-medieval hatha yoga manuals underwent an unrestricted vertical explosion in postural density. Compilations like the \textit{Haṭhayogapradīpikā} and the \textit{Gheraṇḍasaṃhitā} systematically stripped away the material realities of trade routes, espionage, and outcaste survival. By enclosing these rogue physical capabilities within a highly sanitized, respectable theological lexicon, elite text-composers turned the weaponized body inward, repackaging it as a passive daily practice schedule tailored for spiritual liberation and monastic consumption.

\subsection{The Shifting Semantics of Sthāna and Āsana}

This process of theological cleaning relies on a calculated semantic shift across the core technical vocabulary of embodiment. Within early classical layers, the term \textit{sthāna} defined an active, kinetic fighting stance or a strategic military station engineered for territorial defense, while \textit{āsana} denoted the manual contact zone of physical grappling challenges. When captured by late-medieval monastic redactors, these dynamic configurations were aggressively flattened. \textit{Sthāna} was re-coded to signify a static, un-moving monastic hermitage, and \textit{āsana} was reduced to a stationary, decorative physical seat optimized for silent contemplative focus. The dynamic-static continuum was thus permanently severed; the grueling physical maneuvers developed by peripatetic acrobat guilds to endure structural trauma and survive toxic frontier landscapes were frozen into the harmless, respectable configurations of modern monastic catalogs.

\subsection{Fluid Management and the Internalization of Stambha}

The mechanical trajectory of this elite enclosure is explicitly transparent when evaluating the changing parameters of physical locking (\textit{stambhana}). In the ancient strata of the Atharvavedic \textit{Kauśikapaddhati}, \textit{stambha} operated as a literal, veterinary weapon of physical immobilization used to freeze enemy troops or execute \textit{viṣa-stambhana} -- the structural arresting of snake venom within the bloodstream to prevent tissue death. In late-medieval manuals like the \textit{Gheraṇḍasaṃhitā} (3.43--45), this violent external defense mechanism was internalized and re-coded as an energetic breath-retention hold (\textit{stambhakarī mudrā}). Holding the breath for five extended periods was no longer used to isolate physical poison on the battlefield; instead, it was turned inward to serve a new metaphysics of fluid management, designed to arrest the downward flow of generative fluid (\textit{bindu}) and ensure contemplative longevity. The physical container was permanently sealed, transforming an asset of inter-state warfare into an introverted technology of self-cultivation.

\subsection{Postural Proliferation as Spectacular Display Technology}

The final inflation of postural forms is driven not by mystical discovery, but by the requirements of public consumption and state-led soft-power performance. As documented in Jason Birch’s critical edition of the Ujjain manuscript of the \textit{Yogacintāmaṇi} (2024), late-medieval redactors alphabetized eighty-four distinct positions in a structured table, creating the largest compiled grid of physical practice in the premodern record. This scholastic normalization matches the explicit admissions of the 20th-century Marathi physical culture encyclopedia, the \textit{Vyāyām Jñānakośa} (1937--1939). The editors explicitly state that executing complex yoga postures on top of a vertical wooden column (\textit{mallakhāmbāvar āsané}) shares no organic, functional relationship with real combat training; rather, the practice is defined purely as a spectacular display technology (\textit{prekṣaṇīya kām}) engineered to capture civic audiences and project majoritarian national pride.

\section{Conclusion: The Global Topology of Capture}

A rigorous topological analysis of our 20,529-node digital corpus exposes the definitive structural myth of modern ``Yogaland'' and ``Sanskritland'' architectures \cite{McCartney2021, McCartney2026b}. The diachronic network centrality curves document an acute, non-linear structural inversion between independent subaltern mobility and localized somatic enclosure. During the classical Kautilyan and late antique horizons, the \textit{Subversive Mobility Index} reached its historical apex, tracking the intense material reality of clandestine transit, occupational border-evasion, and performative state intelligence infrastructure (\textit{cara}, \textit{sattriṇ}) preserved within the \textit{Arthaśāstra} and narrative layers of the \textit{Kathāsaritsāgara}. Concurrently, the \textit{Somatic Contortion Index} remained dormant on the baseline, verifying that specialized body positioning had not yet been captured or regularized by elite text-composers.

The mid-medieval transition marks the initiation of the structural capture axis. As the volatile, rogue physical capabilities of traveling performer guilds (\textit{ḍomba}, \textit{naṭa}) and outcaste root-cutters (\textit{śvapāka}) were systematically cannibalized by regional courtly networks and early hatha compilers like those of the \textit{Amaraugha} and the Ujjain \textit{Yogacintāmaṇi}, the capacity for independent subaltern mobility was forcefully crushed. This historical sterilization is visually captured by the precipitous crash of the mobility line and the simultaneous, near-vertical explosion of the \textit{Somatic Contortion Index} beyond the 15th century. Reaching an unrestricted apex exceeding 100 points, this modern surge charts the final consolidation of transnational capitalist wellness tourism. The dynamic, dangerous survival tactics of the frontier underbelly are shown to be completely domesticated, turning inward to serve as static, alphabetized, and respectably bourgeois display technologies engineered for public consumption, soft-power branding, and the bio-moral regulation of global civilizational space.

\begin{thebibliography}{99}

\bibitem[Kautilya, 1969]{ArthaŚ}
Kautilya (1969). \textit{Arthaśāstra of Kautilya}. Edited by R. P. Kangle. Bombay: Motilal Banarsidass. (Digital edition integrated via Hellwig's Digital Corpus of Sanskrit node records).

\bibitem[Keśava, 1982]{KauśP}
Keśava (1982). \textit{Kauśikapaddhati of Keśava}. Edited by V. P. Limaye and R. D. Vadekar. Pune: Tilak Maharashtra Vidyapeeth.

\bibitem[Someśvara III, 1939]{Mān_3}
Someśvara III (1939). \textit{Mānasollāsa of King Bhūlōkamalla Someśvara}, Volume II. Edited by G. K. Shorigondekar. Baroda: Oriental Institute (Gaekwad's Oriental Series).

\bibitem[Modha Brahmins, 1964]{MP_7}
Modha Brahmins (1964). \textit{Mallapurāṇa: A Medieval Sanskrit Text on Malla-yuddha}. Edited by B. J. Sandesara and R. N. Mehta. Baroda: Oriental Institute (Gaekwad's Oriental Series).

\bibitem[Kullūkabhaṭṭa, 2003]{GRETIL_vcp}
Kullūkabhaṭṭa (2003). \textit{Manvarthamuktāvalī (Commentary on the Manusmṛti)}. GRETIL Digital Repository.

\bibitem[Rādhākāntadeva, 2010]{DCS_skd}
Rādhākāntadeva (2010). \textit{Śabdakalpadruma Lexicon data layers}. Digital Corpus of Sanskrit (DCS).

\bibitem[Veda Vyāsa, 1969]{HV_App}
Veda Vyāsa (1969). \textit{The Harivaṃśa: Being the Khila or Supplement to the Mahābhārata}. Edited by P. L. Vaidya. Poona: Bhandarkar Oriental Research Institute.

\bibitem[Puranic Composers, 1980]{ViP_5}
Puranic Composers (1980). \textit{The Viṣṇupurāṇa: With the Commentary of Śrīdhara Svāmī}. Edited by J. L. Shastri. Delhi: Motilal Banarsidass.

\bibitem[McCartney, 2026b]{McCartney2026b}
McCartney, Patrick S. D. (2026). \textit{Imagining Sanskritland: `Sanskrit-speaking' Villages, Linguistic Utopias and the Metaphysics of Development}. London and New York: Routledge (Routledge Studies in South Asian History).

\bibitem[McCartney, 2026c]{McCartney2026c}
McCartney, Patrick S. D. (2026). \textit{The Acro-Yoga Complex and Ha\d{t}hayoga's `Contortionist Turn': An Ethno-Historical Study of Ritual and South Asia's `Pole Dancing' Acrobats}. \textit{Abhidha}, 5, Forthcoming.

\bibitem[McCartney, 2026d]{McCartney2026d}
McCartney, Patrick S. D. (2026). \textit{Poison, Vision and Ecstasy: A Transcultural Study of Pharmakon and Ritual Power from the Veda to Greek Mystery Traditions}. \textit{Electronic Journal of Vedic Studies}, Forthcoming.

\bibitem[McCartney, 2021]{McCartney2021}
McCartney, Patrick S. D. (2021). \textit{The Not So United States of Yogaland: Transnational Wellness Capitalism and the Geopolitics of Erasure}. London: Anthem Press.

\bibitem[Birch, 2024]{Birch2024}
Birch, Jason (2024). \textit{The Amaraugha and Amaraughaprabodha of Gorakṣanātha: A Critical Edition with Translation and Study}. Pondicherry: Institut Français de Pondichéry / École française d'Extrême-Orient.

\bibitem[Pinch, 1996]{Pinch1996}
Pinch, William R. (1996). \textit{Peasants and Monks in British India}. Berkeley: University of California Press.

\bibitem[Pinch, 2012]{Pinch2012}
Pinch, William R. (2012). \textit{Warrior Ascetics and Indian Empires}. Cambridge: Cambridge University Press.

\bibitem[Lorenzen, 1978]{Lorenzen1978}
Lorenzen, David N. (1978). Warrior Ascetics in Indian History. \textit{Journal of the American Oriental Society}, 98(1), 61--75.

\bibitem[O'Hanlon, 2007]{OHanlon2007}
O'Hanlon, Rosalind (2007). Military Sports and the History of the Martial Body in India. \textit{Journal of the Economic and Social History of the Orient}, 50(4), 490--523.

\bibitem[Olivelle, 1987]{Olivelle1987}
Olivelle, Patrick (1987). King and Ascetic: State Control of Asceticism in the Artha\v{s}\={a}stra. \textit{Adyar Library Bulletin}, 51, 39--59.

\bibitem[Olivelle, 2011]{Olivelle2011}
Olivelle, Patrick (2011). \textit{Ascetics and Brahmins: Studies in Ideologies and Institutions}. London: Anthem Press.

\bibitem[Clark, 2006]{Clark2006}
Clark, Matthew (2006). \textit{The Daśanāmī-Saṃnyāsīs: The Integration of Ascetic Lineages into an Order}. Leiden: Brill.

\bibitem[Clark, 2018]{Clark2018}
Clark, Matthew (2018). Akhāṛās: Warrior Ascetics. In K. A. Jacobsen, H. Basu, A. Malinar, \& V. Narayanan (Eds.), \textit{Brill’s Encyclopedia of Hinduism Online}. Leiden: Brill.

\bibitem[Danta, 2015]{Danta2015}
Danta, Chris (2015). Acrobats and Ascetics: Peter Sloterdijk and the Aesthetics of Verticality. \textit{European Journal of English Studies}, 19(1), 66--80.

\bibitem[Srinivas, 2012]{Srinivas2012}
Srinivas, Tulasi (2012). Relics of Faith: Fleshly Desires, Ascetic Disciplines and Devotional Affect in the Transnational Sathya Sai Movement. In B. S. Turner (Ed.), \textit{Routledge Handbook of Body Studies} (pp. 185--205). London and New York: Routledge.

\bibitem[Patwardhan, 1927]{Patwardhan1927}
Patwardhan, N. G. (1927). \textit{Malla-Vidy\={a} va Pehalav\={a}n\={i} K\={a}mat [The Science of Wrestling and Pugilistic Maneuvers]}. Poona: Chitrashala Press.

\end{thebibliography}

	
\end{document}
